% Gerado por lcqui-spec-doc; não editar.
\subsection*{M10\_Patrimonio --- formalização executável M10}
Shapes estruturais de bem, resumo catalográfico, requisições de adição/edição, lock, baixa, evento histórico e alteração. CUE verifica campos, enums, nulabilidade e limites locais; identidade Resumo x Bem, máquina de estados, locks, unicidade permanente, versionamento, terminalidade e composição M7/M9 são verificados em Alloy; canonicalização de plaqueta e proveniência por Rust. CUE não prova concorrência temporal.
\begin{description}
\item[id\_resumo\_bem\_patrimonial] Tipo: string. Obrigatório: sim. Aceita nulo: não.
SQL: TEXT.
FK do modelo catalográfico compartilhado.
\item[numero\_patrimonio] Tipo: string. Obrigatório: sim. Aceita nulo: não.
SQL: TEXT.
Plaqueta canônica trim().toUpperCase(), não vazia, no máximo 30 caracteres.
\item[estado\_conservacao] Tipo: enum. Obrigatório: sim. Aceita nulo: não.
SQL: ENUM.
Valores: BOM, REGULAR, RUIM.
\item[id\_local] Tipo: string. Obrigatório: sim. Aceita nulo: não.
SQL: TEXT.
Autoridade de localização; predio/andar/sala são projeções derivadas.
\item[photo\_url] Tipo: string. Obrigatório: sim. Aceita nulo: não.
SQL: TEXT.
Foto obrigatória da unidade física; validação binária pertence ao backend.
\item[documento\_dado\_baixa\_pdf\_url] Tipo: string. Obrigatório: sim. Aceita nulo: sim.
SQL: TEXT.
Nulo antes da baixa; vinculado e imutável no estado terminal.
\item[nome\_responsavel\_sei] Tipo: string. Obrigatório: sim. Aceita nulo: não.
SQL: TEXT.
Responsável institucional; no máximo 150 caracteres.
\item[status] Tipo: enum. Obrigatório: sim. Aceita nulo: não.
SQL: ENUM.
Valores: Ativo, Inservivel, Ja\_dado\_baixa.
\item[descricao\_complementar] Tipo: string. Obrigatório: sim. Aceita nulo: sim.
SQL: TEXT.
Observação da unidade física; no máximo 200 caracteres.
\item[versao] Tipo: int. Obrigatório: sim. Aceita nulo: não.
SQL: INTEGER.
Inicia em 1; incrementa somente em fato canônico, nunca em fan-out derivado.
\item[numero\_patrimonio\_proposto] Tipo: string. Obrigatório: sim. Aceita nulo: não.
SQL: TEXT.
Plaqueta bruta informada na requisição de adição.
\item[numero\_patrimonio\_normalizado] Tipo: string. Obrigatório: sim. Aceita nulo: não.
SQL: TEXT.
N(proposto) = trim().toUpperCase(); única chave de lock e Chaves\_Unicas.
\item[versao\_bem\_origem] Tipo: int. Obrigatório: sim. Aceita nulo: não.
SQL: INTEGER.
Versão observada para controle otimista de edição.
\item[id\_requisicao] Tipo: string. Obrigatório: sim. Aceita nulo: não.
SQL: TEXT.
Requisição proprietária do lock; precisa coincidir antes da liberação.
\item[tipo\_lock] Tipo: enum. Obrigatório: sim. Aceita nulo: não.
SQL: ENUM.
Valores: EDICAO, ADICAO.
\item[chave\_recurso] Tipo: string. Obrigatório: sim. Aceita nulo: não.
SQL: TEXT.
ID do bem (edição) ou plaqueta normalizada (adição).
\item[tipo\_historico] Tipo: enum. Obrigatório: sim. Aceita nulo: não.
SQL: ENUM.
Valores: cadastro, edicao, baixa.
\end{description}
